\section{Zielbild der SARIA IT}\label{sec:5zielbild}
Im folgenden Kapitel wird die strategische Ausrichtung der SARIA IT näher dargestellt und in Bezug auf die Analyse der aktuellen Situation ein Soll-Zustand beschrieben
\subsection{Zielbild der internationalen IT}

Im Zuge der Professionalisierung der deutschen IT und dem UNO-Projekt\footnote{siehe Kapitel \ref{subsubsec:3projectuno}}, wurde eine Strategie für die Zukunft der SARIA IT festgelegt.

Als Teil dieser wurde eine Vision beschrieben, die als Ziel für die nächsten Jahre gelten soll.
\glqq Becoming SARIA's common and professional Group IT\grqq{} beschreibt das Ziel, die SARIA IT auf einem globalen Level (vorest im Scope von UNO\footnote{siehe Kapitel \ref{subsubsec:3projectuno} auf S. \pageref{subsubsec:3projectuno}}) eine gemeinsame und professionelle IT zu etablieren.

Die Mission der SARIA Group IT lautet: \glqq Our mission is to enable the Divisions and Corporates to work across countries and focus on the business. We achieve this by growing together into a strong Group IT and by becoming a professional and internal service partner.\grqq{}
\\
Diese beschreibt den Zweck der IT treffend, in dem diese sich als Enabler darstellt, welche die globale Zusammenarbeit ermöglichen kann und dabei das Business in den Mittelpunkt ihrer Aktivitäten stellt.
Dazu will die IT zusammenwachsen, aus den lokalen Einheiten heraus und ein professioneller Service-Partner werden.
\\
Weiterhin hat die IT in der Strategie 4 Fokus-Bereiche festgelegt:
\begin{figure} [H]
	\centering
	\includegraphics[page={14},scale=0.4, keepaspectratio]{resources/images/IT SUMMIT 2023 - final.pdf}
	\caption{Fokus-Bereiche der IT-Strategie für die nächsten Jahre}
	\label{fig:anforderungenitsm}
\end{figure}

Es soll eine gemeinsame Group IT kreirt werden, die Services harmonisert und gemeinsame Synergien nutzen kann.
Zudem soll die Expertise der einzelnen Personen erhöht werden und ein höherer Reifegrad als bisher erreicht werden. Dies soll insbesondere durch die Implementierung von Prozessen auch im Bereich \ac{itsm} erfolgen.
Ein weiterer Schritt soll die Einführung von \ac{sla}s bringen, indem diese die Serviceleistungen der Group IT transparenter darstellen.

Der zweite Fokus Bereich der Strategie zielt auf die Harmonisierung der Kommunikation und Kollaboration ab.
Auch hier sollen Synergien verwendet werden, insbesondere der Umzug der einzelnen Länder auf eine gemeinschaftlich genutzte Plattform steht hier im Vordergrund.
Dies soll es den Kunden der IT ermöglichen, über die Ländergrenzen hinweg einfach und schnell kollaborieren zu können.
Gleichzeitig sieht die Strategie vor, dass so die Möglichkeiten und Services für alle Kunden standardisiert werden um diese auf dem gleichen Service Level anzubieten zu können.

Ein weitere Fokus Bereich verspricht Verbesserungen im Bereich der Datenhaltung.
Hier sollen unterschiedliche Data Environments miteinander verbunden werden. Auch eine globale Datenstrategie wird angestrebt.
Mit diesen Umstrukturierungen sollen auch organisatorische Änderungen umgesetzt werden, um professionelle Prozesse zu ermöglichen und eine bessere Serviceleistung zu erreichen.
\\
Der letzte Fokus-Bereich konzentriert sich auf IT-Sicherheit und die Geschäftsfortführung, insbesondere bei Ausfällen, Angriffen, \acs{etc}.
Dazu soll ein IT-Security Team etabliert werden, welches auch die notwendigen Prozesse definieren kann.
Angestrebt werden moderne und state-of-the-art Sicherheitsmechanismen. Die Errichtung oder Auslagerung eines Managed Security Operation Centers ist ein weiterer wichtiger Teil der Implementierung der IT-Security.
Um die Geschäftsfortführungsprozesse weiter zu verbessern, sollen weitere Maßnahmen ergriffen werden.

\subsection{Zielbild des ITSM}\label{subsec:5}
Die Erfüllung der oben benannten Strategie, soll insbesondere durch organisatorische und prozessuale Veränderungen erreicht werden.
Dazu soll auch die \ac{itsm}-Lösung beitragen.
Mit der Einführung von \ac{itsm}-Prozessen strebt die SARIA IT einen höheren Reifegrad an.
Die bisherigen Versuche ohne passendes Framework im Hintergrund die Verbesserung der Servicequalität zu erreichen, schlugen insbesondere aufgrund von Fehlern in der tatsächlichen Umsetzung fehl.
Daher hat die SARIA IT festgelegt mit Hilfe von ITIL als Grundlage einige Prozesse einzuführen und so zu professionalisieren.
Dabei soll \ac{itil} nur als Unterstützung bieten, eine vollkommene Umsetzung wird zum aktuellen Zeitpunkt nicht in Erwähgung gezogen und würde auch dem Ansatz von \ac{itil} widersprechen.